\section{Fragenpool erstellen und verwalten}
\label{chap:fragenpool_erstellen}
\setauthor{\secondauthor}

\begin{figure}[H]
\centering
\fbox{
\includegraphics[width=0.7\textwidth]{FragenpoolHinzufügen.png}
}
\caption{Zentrale Oberfläche zur individuellen Zusammenstellung des Fragenpools}
\label{fig:fragen-browsen}
\end{figure}

Bevor Nutzer:innen mit dem eigentlichen Lernen starten können, ist das Erstellen eines persönlichen Fragenpools nötig. Wenn Nutzer:innen im Partnerprojekt bereits selbst Fragen erstellt haben, wurden diese automatisch in den Fragenpool übernommen.

\subsection{Die „Fragen browsen“-Ansicht als digitaler Katalog}
\label{chap:fragen_browsen}
\setauthor{\secondauthor}

Die „Fragen browsen“-Ansicht fungiert als Herzstück der Inhaltsverwaltung. Sie ist als interaktiver, digitaler Katalog konzipiert, der den Zugriff auf die umfangreiche Fragenauswahl der Plattform ermöglicht. Grundlage bildet eine klare Trennung zwischen der Navigation und dem Hauptbereich mit den Fragen.

Das Interface ist für eine intuitive Bedienung in zwei funktionale Hauptbereiche unterteilt:

\begin{itemize}
\item \textbf{Strukturierte Navigations-Seitenleiste (links)}:  
Um die Komplexität der verschiedenen Lehrpläne abzubilden, nutzt die Seitenleiste eine hierarchische Baumstruktur. Nutzer:innen finden hier eine geordnete Auflistung aller verfügbaren Hauptfächer. Durch ein interaktives Dropdown-System lassen sich diese Fächer in ihre jeweiligen Themengebiete und Unterkategorien aufgliedern. Diese Struktur erlaubt es, innerhalb weniger Klicks von einer globalen Fachansicht (z.\,B. Mathematik) zu einer spezifischen Problemstellung (z.\,B. Differenzialrechnung) zu navigieren. Die optische Hervorhebung des aktuell gewählten Pfades dient dabei als Orientierungshilfe.

\item \textbf{Hauptbereich mit Vorschaufunktion (rechts)}:  
Im Fokus des Hauptfensters stehen die Fragenkarten. Jede Karte ist so gestaltet, dass sie bereits in der Kompaktansicht die wesentlichen Informationen vermittelt. Ein wesentliches Merkmal ist die integrierte \textbf{Schnellvorschau}. Erst beim Ausklappen einer Fragekarte werden alle Details sowie ein „Üben“-Button, das zugehörige Bild und der vollständige Fragentext angezeigt. So können Nutzer:innen eine Frage probeweise üben, ohne die Verwaltungsansicht verlassen zu müssen und der Lernprozess somit nicht gestört wird. Dadurch lässt sich direkt beurteilen, ob Schwierigkeitsgrad und Inhalt zur aktuellen Lernphase passen.
\end{itemize}

\begin{figure}[H]
\centering
\fbox{
\includegraphics[width=0.9\textwidth]{FragenBrowsenÜben.png}
}
\caption{Oberfläche zum direkten Testen einer Frage aus der Browsing-Ansicht}
\label{fig:fragen-browsen-uben}
\end{figure}

\subsection{Auswahl und individuelle Anpassung des Fragenpools}
\label{chap:auswahl_logik}

\textbf{Dynamisches Hinzufügen}

Das bekannte \textbf{Stern-System} wird hier verwendet um zu symbolisieren das eine Frage im Fragenpool ist. Angelehnt an bekannte Favoriten-Funktionen aus modernen Webanwendungen können Nutzer:innen durch Anklicken des Sternsymbols eine Frage direkt ihrem aktiven Lernpool hinzufügen. Die Auswahl wird sofort visuell durch eine Farbänderung des Symbols bestätigt. Dieser Vorgang ist bewusst einfach und unkompliziert gestaltet, um den Lernfluss beim Durchsehen der Kataloge nicht zu unterbrechen. Ebenso leicht lassen sich Fragen wieder entfernen, so bleibt der Pool jederzeit flexibel anpassbar.

\textbf{Ansichtssteuerung: „Alle“ / „Mein Pool“}

Um auch bei hunderten verfügbaren Fragen den Überblick zu behalten, verfügt die Ansicht über einen umschaltbaren Filtermodus. Dieser unterstützt die Selbstorganisation beim Lernen:

\begin{itemize}
\item \textbf{Explorations-Modus („Alle“)}: In dieser Ansicht werden die gesamten verfügbaren Fragen sichtbar. Sie dient dem Entdecken neuer Inhalte und zeigt auch jene Fragen an, die noch nicht Teil des persönlichen Fragenpools sind.
\item \textbf{Fokus-Modus („Bereits im Fragenpool“)}: Diese Ansicht funktioniert wie ein digitaler Schreibtisch, auf dem nur die aktuell relevanten Fragen liegen. Nutzer:innen können hier schnell überprüfen, welche Fragen bereits ausgewählt wurden, und in welche Themenbereiche noch im individuellen Fragenpool fehlen.
\end{itemize}

\textbf{Effizienzsteigerung durch Batch-Operationen}

Da die Vorbereitung auf die Matura oft unter Zeitdruck erfolgt, bietet das System praktische Funktionen zur Massenverarbeitung. Anstatt jede der teils über 50 Fragen eines Kapitels einzeln auswählen zu müssen, erlaubt der \textbf{„Alle auswählen“-Button} die schnelle Übernahme ganzer Module. Dies ist besonders hilfreich für Nutzer:innen, die sich zu Beginn eines Lernzyklus ein komplettes Fachgebiet erschließen möchten. In Kombination mit der anschließenden gezielten Abwahl einzelner Fragen ermöglicht dieser „Top-Down-Ansatz“ eine sehr effiziente und dennoch präzise Zusammenstellung des persönlichen Lernumfelds.

Durch diese flexible Form der Verwaltung wird gewährleistet, dass nachfolgende Module – wie die Startansicht (Kapitel \ref{chap:startansicht}) und der Übungsmodus (Kapitel \ref{chap:uebung}) – auf einer individuell abgestimmten Datenbasis aufbauen können. Nutzer:innen behalten dabei jederzeit die Kontrolle über ihre Lerninhalte und können diese aktiv gestalten.